\chapter{Adaptive Interpreters and Staged Compilation}
\label{ch:adaptive-staged}

An addition site may execute millions of times with integers and once with strings.
A generic interpreter can handle both, but it repeatedly pays for cases that rarely occur.
Compiling the whole application before seeing any inputs spends effort on code that may never run.
Adaptive execution starts with a general implementation, specializes where observations justify the expense, and retains a correct route for later surprises.

Adaptation and staging solve related but different problems.
Adaptation changes execution strategy as runtime observations arrive.
Staging separates computations according to which inputs are already known.
A self-optimizing interpreter can provide a good substrate for a partial evaluator, but neither idea requires the other.

\section{Quicken one operation before compiling a function}
\label{sec:ch65-quickening}

\term{Quickening} replaces a general execution action with a more specialized action after information becomes available.
The replacement may be a bytecode opcode, an AST node, or an entry in a side table.
The physical representation matters less than the dispatch contract: the fast action must check or otherwise establish its assumptions.

Our addition site accepts exactly two integers or exactly two strings.
Booleans are excluded even though Python treats its Boolean type as an integer subclass.
This explicit choice avoids accidentally inheriting host-language semantics in the guest language.
The site begins cold, changes to an integer-specialized state after an integer pair, and moves permanently to generic after a mismatch.
Algorithm~\ref{alg:ch65-quicken} performs no observable action before the guard.

\begin{algorithm}[htbp]
\caption{Execute a guarded addition site with one-way generalization.}
\label{alg:ch65-quicken}
\begin{pseudocode}
\Require operands $a,b$, site state in cold, integer, or generic
\Ensure the language's integer/string addition result or type error
\State $ints\gets$ both operands have exactly the integer type
\If{the site is cold}
  \State select integer state if $ints$, otherwise generic
\EndIf
\If{the site is integer and $ints$ is false}
  \State increment miss count; replace the site state with generic
\EndIf
\If{$ints$}
  \State \Return integer addition
\EndIf
\If{both operands have exactly the string type}
  \State \Return string concatenation
\EndIf
\State \Return type error
\end{pseudocode}
\end{algorithm}

The state transition and type checks take constant time; integer arithmetic and string construction cost depends on operand size.
The Python model makes the state transitions visible but does not claim a dispatch-speed improvement: it still uses host-language conditionals to demonstrate the contract.
A bytecode interpreter can install a different handler, while an AST interpreter can replace the executing node~\cite{wurthinger-2012-interpreters}.

Table~\ref{tab:ch65-site} works three calls through the site.
Returning to integers does not re-specialize in this model.
That policy avoids oscillation but may sacrifice later performance; a production system can use counters, cooldowns, or polymorphic specializations.

\begin{table}[htbp]
\caption{A mismatch widens the site's behavior before evaluating the operation.}
\label{tab:ch65-site}
\centering
\begin{tabular}{llll}
\toprule
Operands & Previous state & Result & New state \\
\midrule
$3,4$ & cold & $7$ & integer \\
\code{a}, \code{b} strings & integer & \code{ab} & generic \\
$8,9$ & generic & $17$ & generic \\
\bottomrule
\end{tabular}
\end{table}

Figure~\ref{fig:ch65-site} presents the state graph.
An exception is an outcome of executing the operation, not a fourth specialization state.
In a multi-threaded interpreter, replacing a handler or AST node also requires a publication protocol so readers never observe partially initialized specialization data.

\begin{figure}[htbp]
\centering
\begin{tikzpicture}[node distance=16mm and 20mm]
\node[basic block] (cold) {cold};
\node[basic block,below left=of cold] (ints) {integer};
\node[basic block,below right=of cold] (generic) {generic};
\draw[cfg edge] (cold)--node[left,font=\small] {integer pair}(ints);
\draw[cfg edge] (cold)--node[right,font=\small] {otherwise}(generic);
\draw[cfg edge] (ints)--node[below,font=\small] {mismatch}(generic);
\end{tikzpicture}
\caption{One-way generalization bounds the number of state changes at the addition site.}
\label{fig:ch65-site}
\end{figure}

\section{Fuse dispatch while preserving observation points}
\label{sec:ch65-superinstructions}

A \term{superinstruction} combines a common sequence of interpreter actions into one dispatch unit.
Loading two locals and adding them may become a single handler.
This reduces dispatch overhead without necessarily changing the arithmetic representation.
It differs from specializing the addition for integers, although the two techniques can be combined.

Fusion must preserve the guest stack effect, exceptions, and any required intermediate observation points.
If the second load can fail, the fused handler must report the corresponding guest location.
If a debugger can stop between the loads and the addition, the runtime needs a way to reconstruct that intermediate state or a policy defining which boundaries remain observable.
Signal, garbage-collection, and tracing requirements are part of this contract, not afterthoughts to opcode selection.

Demand-driven compilation similarly requires a stable entry mechanism.
An unresolved entry can invoke a compiler, install the resulting code, and then transfer control with the original arguments.
Concurrent first calls need a materialization state so they do not publish incompatible duplicate definitions.
LLVM ORC explicitly models symbol materialization and supports lazy compilation arrangements~\cite{llvm-2026-orc}.

\section{Specialize an interpreter against a program}
\label{sec:ch65-staging}

Let $I(p,x)$ denote an interpreter running program $p$ on input $x$.
A \term{partial evaluator} receives a computation and some known inputs and produces a residual computation for the remaining inputs.
If $p$ is known and $x$ is dynamic, the desired equation is
\[
I(p,x)=I_p(x),
\]
where $I_p$ is the residual program obtained by specializing the interpreter with respect to $p$.
Equality here means the specified observable behavior, including errors and termination, rather than just a successful numeric result.

Our executable staging example uses expression syntax containing constants, one dynamic input, addition, and multiplication over mathematical integers.
The specializer traverses the known syntax once and emits stack bytecode.
It removes runtime traversal of the expression tree, but does not perform constant folding or generate native code.
It is a hand-written specialization for this interpreter, not a general self-applicable partial evaluator.

\begin{algorithm}[htbp]
\caption{Residualize known expression syntax into stack bytecode.}
\label{alg:ch65-stage}
\begin{pseudocode}
\Require finite expression tree with constants, input, addition, and multiplication
\Ensure bytecode with one result and the same integer expression meaning
\Function{Emit}{$e$}
  \If{$e$ is a constant}
    \State append a constant-push instruction
  \ElsIf{$e$ is the dynamic input}
    \State append an input-push instruction
  \Else
    \State \Call{Emit}{left child of $e$}; \Call{Emit}{right child of $e$}
    \State append the corresponding binary operation
  \EndIf
\EndFunction
\end{pseudocode}
\end{algorithm}

Algorithm~\ref{alg:ch65-stage} uses a shared output buffer, giving $O(n)$ construction time and output space for $n$ syntax nodes, with recursion depth equal to tree height.
The compact Python companion concatenates lists instead, so its worst-case construction cost is $O(n^2)$ on a skewed tree.
This difference is an implementation detail worth measuring rather than hiding behind the linear conceptual traversal.

For $(x+5)\cdot3$, the residual sequence pushes input, pushes $5$, adds, pushes $3$, and multiplies.
With $x=4$, stack states are $[4]$, $[4,5]$, $[9]$, $[9,3]$, and $[27]$.
Figure~\ref{fig:ch65-staging} separates compile-time syntax from runtime input.
The tests compare source interpretation and residual execution over $17$ constant choices and $201$ input values.

\begin{figure}[htbp]
\centering
\begin{tikzpicture}[node distance=14mm and 16mm]
\node[basic block] (syntax) {known expression syntax};
\node[basic block,below=of syntax] (code) {residual stack bytecode};
\node[basic block,right=of code] (input) {dynamic input x};
\node[basic block,below=of code] (result) {execution result};
\draw[cfg edge] (syntax)--node[right,font=\small] {specialize}(code);
\draw[cfg edge] (code)--(result); \draw[cfg edge] (input)--(result);
\end{tikzpicture}
\caption{Staging consumes the known program structure before dynamic inputs arrive.}
\label{fig:ch65-staging}
\end{figure}

\section{Understand the Futamura projections precisely}
\label{sec:ch65-projections}

Write $S(f,a)$ for specializing a two-input program $f$ with its first input fixed to $a$.
With compatible encodings and a suitable self-applicable specializer, the first projection is $S(I,p)$: a compiled form of $p$.
The second is $S(S,I)$: a compiler that maps programs to those residual forms.
The third is $S(S,S)$: a compiler generator that takes an interpreter and produces a compiler~\cite{futamura-1999-partial}.
Table~\ref{tab:ch65-projections} records the inputs that remain dynamic at each level.

\begin{table}[htbp]
\caption{Each projection fixes another layer of interpretive machinery.}
\label{tab:ch65-projections}
\centering
\begin{tabular}{lll}
\toprule
Expression & Produced artifact & Remaining input \\
\midrule
$S(I,p)$ & residual program & program data \\
$S(S,I)$ & compiler & source program \\
$S(S,S)$ & compiler generator & interpreter \\
\bottomrule
\end{tabular}
\end{table}

These equations express semantic relationships, not guarantees of useful code quality or feasible specialization time.
Binding-time analysis decides which computations may run early.
Code growth, static recursion, dynamic control flow, and effect duplication can make an apparently simple specializer unusable.
An effectful input read must remain residual unless the staging contract explicitly makes that read static.

\section{Generate interpreters and bootstrap metacircular systems}
\label{sec:ch65-generated}

An opcode specification can drive several generated consumers.
For a stack addition, the specification might record two pops, one push, integer operand kinds, an arithmetic operation, and a source-location rule for errors.
A generator can use those fields to emit a decoder table, a verifier transfer, and an interpreter-handler skeleton.
The generated components then share numeric tags and operand shapes rather than duplicating them by hand.
The generator still needs a semantic definition for the handler body; a stack-effect tuple alone cannot distinguish addition from subtraction.

Consider generating the residual machine used here.
The input instruction pushes one dynamic value; a constant instruction pushes its payload; addition and multiplication consume two values and push one.
Those shapes suffice to generate structural stack checks.
They do not establish that the implementation pops the operands in the intended order, uses arbitrary-precision arithmetic, or reports the correct source continuation on failure.
If subtraction were added, a reversed pop order would preserve every stack depth while computing the wrong answer.
Differential tests against the expression interpreter therefore remain useful even when both consumers originate from one schema.

A \term{metacircular interpreter} is written in the language it interprets, or in a closely corresponding language subset under a stated definition.
That relationship can clarify semantics and support reuse, but does not make the interpreter a self-applicable specializer.
Likewise, generating an interpreter from an opcode specification is code generation, not automatically the second Futamura projection.
The distinctions keep useful implementation techniques from being credited with capabilities they do not implement.

A metacircular system needs an initial execution path: a trusted host implementation, a small bootstrap interpreter, or an existing compiled version.
Running the language's interpreter through itself adds another interpretive layer unless specialization or compilation removes it.
Features delegated to the host must be identified, especially arithmetic overflow, object identity, exceptions, and evaluation order.
The exact-type checks in our quickening example illustrate the issue: adopting Python's Boolean/integer subtype relationship accidentally would change the guest language before any optimization occurred.

For a practical bootstrap, first compare the host-executed interpreter with the interpreter running guest programs through the bootstrap layer on small cases.
Then validate any generated or specialized implementation against the same observable contract.
This establishes a chain of explicit comparisons rather than assuming that using one language to implement itself confers correctness.
The companion implements the expression/residual comparison only; a schema generator and a self-interpreting language are extensions with their own tests.

\section{Exercises}
\label{sec:ch65-exercises}

\begin{enumerate}
\item\label{ex:ch65-transition} Run the site on $(3,4)$, two strings, $(8,9)$, and a Boolean/integer pair.
Give states, miss count, results, and the final error.
\item\label{ex:ch65-stack} Emit residual bytecode for $(x+2)\cdot4$ and show its stack trace at $x=3$.
\item\label{ex:ch65-effect} An interpreter reads a file as part of evaluating a source expression.
Explain why knowing the syntax does not necessarily permit reading the file during specialization.
\item\label{ex:ch65-projection} Which projection produces a compiler from an interpreter, and what self-application capability does it require?
Explain why the companion does not implement that projection.
\end{enumerate}

\section{Further reading}
\label{sec:ch65-reading}

Self-optimizing AST interpreters provide a concrete approach to runtime node specialization~\cite{wurthinger-2012-interpreters}.
Futamura's article develops the specialization relationships summarized above~\cite{futamura-1999-partial}.
ORC's documentation supplies an implementation-oriented account of deferred materialization~\cite{llvm-2026-orc}.
